La ley\par

\[
M_n
=
\left(\frac{a_n}{a_0}\right)^{-6}
\]

is perhaps the simplest algebraically and one of the deepest conceptually.\par

The Perron dynamics produces two complementary behaviors.\par

On the one hand, there exists an expansion:\par

\[
V_n=\varphi^n.
\]

In three spatial dimensions, the associated linear scale is the cube root:\par

\[
\frac{a_n}{a_0}
=
\varphi^{n/3}.
\]

On the other hand, there exists a stable mode of memory:\par

\[
M_n=\varphi^{-2n}.
\]

Expansion grows with \(\varphi^n\). Stable memory decreases with \(\varphi^{-2n}\). When level \(n\) is eliminated, there necessarily appears\par

\[
M_n
=
\left(\frac{a_n}{a_0}\right)^{-6}.
\]

The exponent \(-6\) is derived from two internal facts:\par

\begin{itemize}[leftmargin=*,itemsep=0.35em]
\item the space has dimension \(3\), so the linear scale carries exponent \(1/3\);
\item the stable mode is quadratic and therefore has exponent \(-2\).
\end{itemize}

The reason between both is\par

\[
\frac{-2}{1/3}=-6.
\]

In Einstein–Cartan, the spin density scales as an inverse volume:\par

\[
s\sim a^{-3}.
\]

Its quadratic contribution then scales as\par

\[
s^2\sim a^{-6}.
\]

The HMT law reproduces that exponent exactly before introducing the Einstein–Cartan realization. This is the genuinely suggestive feature.\par

The Perron memory mode behaves as a quadratic spin density should behave in a geometry with torsion.\par

The natural interpretation is that torsion physically publishes the genealogical memory complementary to metric information.\par

Curvature describes how geometry deforms. Torsion describes how geometry retains a transport orientation in addition to distances. HMT possesses orientation, carry, and route memory from the outset. The correspondence with Einstein–Cartan is therefore structural: both frameworks distribute geometric history between metric and torsion.\par

When the universe expands, that memory is diluted as \(a^{-6}\). It is diluted faster than an ordinary matter or radiation density. But precisely for that reason it may is dominant in the small-scale region.\par

In a primordial regime, a torsional contribution may compete with ordinary density and modify the singularity. The possibility of a bounce depends on the sign, amplitude, spin current, and boundary conditions. HMT first provides the exact scaling law; the physical realization then fixes those conditions.\par

That offers a very clear division of labor:\par

\begin{itemize}[leftmargin=*,itemsep=0.35em]
\item HMT generates the stable memory and its exponent;
\item the Einstein–Cartan realization assigns that memory to a quadratic spin density;
\item cosmological dynamics fixes amplitude, sign, and physical conditions.
\end{itemize}

The relation to Mach's principle appears when this law is combined with the derivation of \(G\).\par

In the gravitational derivation, the local coupling is selected through the closure of a global cycle. In the cosmological law, the local torsion state preserves the memory of global expansion. In both cases, local physics is defined relationally by the global genealogy.\par

The local geometry recalls the global state.\par

This is the most concrete sense in which the construction is machina:\par

\begin{itemize}[leftmargin=*,itemsep=0.35em]
\item inertia participates in the global clock;
\item the gravedad belongs to the clausura of phase;
\item torsion embodies the memory of history;
\item the local scale encodes the global expansion level.
\end{itemize}

Mach's principle then acquires an operatorial form: local magnitudes are evaluations of a global genealogy, and physical invariants are what survives transport between the two scales.\par

The branch of temporal periodicity of (108) phases and the Perron branch also express two complementary aspects of cosmology.\par

The temporal periodicity of (108) phases fixes to periodicity scale and, through to of Sitter realization, to curvature scale. It is the cosmology of closure: it asks which radius allows the cycle to is coherent.\par

Perron fixes an expansion rate and a law of memory dilution. It is the cosmology of evolution: it asks how the genealogy changes as levels are traversed.\par

One branch describes the closure of geometry. The other describes the persistence of its memory.\par

The chronological calendar and Perron dynamics remain typologically distinct and complementary realizations. Together they produce an especially rich image: the universe possesses a global phase scale and an internal dynamics of memory. Curvature and torsion are two publications of those two aspects.\par

\section{Confluence of vacuum readers, action, gravity, area, mixture, and memory}

If the nine results are read together, a common grammar appears.\par

Each reader selectively publishes an invariant and carries the rest as complementary memory.\par

The vacuum preserves:\par

\begin{itemize}[leftmargin=*,itemsep=0.35em]
\item the product as propagation;
\item the quotient as orientation.
\end{itemize}

The pair \(G,\hbar\) conserves:\par

\begin{itemize}[leftmargin=*,itemsep=0.35em]
\item action as orientation;
\item area as the common geometry.
\end{itemize}

The gravitational closure preserves:\par

\begin{itemize}[leftmargin=*,itemsep=0.35em]
\item the equality among phase, Compton, gravity, and Planck.
\end{itemize}

The mode \(30\) preserves:\par

\begin{itemize}[leftmargin=*,itemsep=0.35em]
\item the common cadence of vacuum, criticality, and memory.
\end{itemize}

The modular area and Hamiltonian are conserved:\par

\begin{itemize}[leftmargin=*,itemsep=0.35em]
\item the total quantity;
\item sectoral provenance.
\end{itemize}

The transito \(6\to8\) preserves:\par

\begin{itemize}[leftmargin=*,itemsep=0.35em]
\item the geometric turn in Barbero;
\item its informational footprint in entropy;
\item its mixing orientation in CKM.
\end{itemize}

The angular determinant preserves:\par

\begin{itemize}[leftmargin=*,itemsep=0.35em]
\item the even part that can become electroweak mixing.
\end{itemize}

The quarter-turn torsion preserves:\par

\begin{itemize}[leftmargin=*,itemsep=0.35em]
\item a complex structure published as an arithmetic character.
\end{itemize}

The Perron dynamics preserves:\par

\begin{itemize}[leftmargin=*,itemsep=0.35em]
\item the stable memory of the expansion, realized as cosmological torsion.
\end{itemize}

For that reason the constants are related through to common architecture of transport and preserve their functional identities. They are different residues of that architecture.\par

Triadic Modular Holography provides the genealogy.\par
Dimensional Mechanics provides the realization.\par
The real number appears finally as the terminal value.\par

In that view, a fundamental constant is a complete statement about the universe. It says:\par

\begin{itemize}[leftmargin=*,itemsep=0.35em]
\item what changed;
\item which orientation was inverted;
\item which memory crossed the boundary;
\item what information remained interior and what part was published;
\item which combination remained invariant;
\item and the physical magnitude in which that invariance was ultimately realized.
\end{itemize}

That is what makes these derivations remarkable. They jointly explain the values of the constants and why each one performs exactly its structural function.\par
